% Chapter 7: Graham 2001 - Positivity in Equivariant Schubert Calculus
% Based on William Graham, Duke Mathematical Journal 109(3):599-614, 2001
% arXiv: math/9908172

\chapter{Graham's Theorem: Positivity in Equivariant Schubert Calculus}
\label{chap:GrahamPositivity}

\section{Introduction and Historical Background}

In 2001, William Graham published a paper in the Duke Mathematical Journal \cite{Gr} that proved a central conjecture in equivariant Schubert calculus---the Dale Peterson conjecture---establishing the positivity property of coefficients in the equivariant cup product expansion. This result represents a milestone in the development of Schubert calculus.

\subsection{Origin of the Problem}

In classical Schubert calculus, given three Schubert cells $X_u, X_v, X_w$ on a flag manifold $G/P$, their intersection number is given by the Schubert structure constants $c_{u,v}^w$. Since geometric intersections are transverse in general position, $c_{u,v}^w$ is a well-defined nonnegative integer. This forms the foundation of classical Schubert calculus.

However, when one considers equivariant cohomology $H_T^*(G/P)$, the problem becomes more refined: the action of the torus $T$ introduces additional weight parameters. The product expansion coefficients $c_{u,v}^w(t,s)$ are no longer merely integers, but polynomials in the torus weight differences $t_i - s_j$. The question is: do these coefficients still maintain positivity?

\subsection{The Peterson Conjecture}

Dale Peterson proposed a bold conjecture: in equivariant Schubert calculus, when $u, v$ are Grassmannian permutations, the expansion coefficients $c_{u,v}^w(t,s)$ are polynomials with nonnegative coefficients in $t_i - s_j$.

This conjecture carries profound significance in both algebraic geometry and combinatorics:
\begin{itemize}
\item It generalizes the classical result that $c_{u,v}^w \in \mathbb{Z}_{\geq 0}$
\item It reveals deep connections between equivariant cohomology and the geometry of flag manifolds
\item It laid the foundation for the later quantum version (Mihalcea's work)
\end{itemize}

Graham's paper gave the first complete proof of this conjecture.

\section{Main Theorem}

\begin{Theorem}[\cite{Gr} Theorem 1.3]
\label{th:GrahamMainPositivity}
Let $u, v \in S_n$ be Grassmannian permutations. Then in the equivariant cohomology $H_T^*(Fl_n)$ of the flag manifold $Fl_n$, the cup product of the Schubert classes $\sigma_u, \sigma_v$ satisfies
\begin{equation}
\sigma_u \cdot \sigma_v = \sum_{w \in S_n} c_{u,v}^w(t,s) \sigma_w,
\label{eq:GrahamPositivity}
\end{equation}
where the structure constants $c_{u,v}^w(t,s)$ are polynomials with nonnegative coefficients in $t_i - s_j$ ($1 \leq i, j \leq n$), namely
$$c_{u,v}^w(t,s) \in \mathbb{N}[t_i - s_j]_{i,j \geq 1}.$$
\end{Theorem}

\begin{Remark}
Here $t_i$ and $s_j$ are formal variables for the torus weights in equivariant cohomology. When we set all $t_i = s_j = 0$, \cref{eq:GrahamPositivity} reduces to the classical Schubert calculus formula
$$\sigma_u \cdot \sigma_v = \sum_w c_{u,v}^w \sigma_w,$$
where $c_{u,v}^w \in \mathbb{Z}_{\geq 0}$ are the well-known nonnegative integers.
\end{Remark}

\section{Affine Schubert Varieties and the Geometric Framework}

The central innovation in Graham's proof is the introduction of affine Schubert varieties and the establishment of deep connections between them and the quantum homology of flag manifolds.

\subsection{The Affine Grassmannian}

Let $G = \mathrm{GL}_n(\mathbb{C})$, and consider the affine Grassmannian
$$Gr = G(\!(t)\!)/G[[t]],$$
which is an infinite-dimensional algebraic variety. The affine Grassmannian plays a central role in representation theory and geometry; its cohomology is closely related to the quantum cohomology of flag manifolds.

The affine Grassmannian $Gr$ can be viewed as the moduli space of all crystallizable closed subspaces of complex $n$-dimensional vector spaces.\footnote{This perspective connects to the theory of affine Lie algebras and their representations; see \cite{Gr} for details.}

\subsection{Affine Schubert Varieties}

In the affine Grassmannian $Gr$, elements $w$ of the affine Weyl group $\widetilde{W}$ correspond to affine Schubert varieties $\mathcal{X}(w)$. These varieties satisfy the following inclusion relations:
$$\overline{\mathcal{X}(u)} \subseteq \overline{\mathcal{X}(v)} \iff u \leq v \quad (\text{Bruhat order}).$$

Graham's key observation is that affine Schubert varieties admit a cell structure compatible with the moment map, which allows one to establish geometric connections with ordinary flag manifolds.\footnote{The compatibility with the moment map is crucial for establishing the transversality results needed in the proof.}

\section{Peterson Degeneration}

The central tool in the proof is the Peterson degeneration of the quantum homology ring.

\subsection{The Peterson Isomorphism}

Let $QH^*(Fl_n)$ denote the (small) quantum homology of the flag manifold. Peterson proved that there exists a ring isomorphism
$$QH^*(Fl_n) \cong H^*(Gr) \otimes \mathbb{Q}[q] / (q_1, \ldots, q_{n-1}),$$
where $Gr$ is the affine Grassmannian and $q_i$ are the quantum parameters.\footnote{The quantum parameters $q_i$ correspond to the Chern classes of certain line bundles on the flag manifold; their role in the degeneration is explained in \cite{Gr} Section 2.}

This isomorphism degenerates quantum multiplication to ordinary cohomology multiplication together with a limit over a flat family.

\subsection{Degeneration to the Affine Grassmannian}

Graham extended this idea to the equivariant setting, establishing a degeneration between equivariant quantum homology and equivariant cohomology of the affine Grassmannian. Specifically, there exists a flat family such that as the parameter tends to infinity, the equivariant Schubert classes tend to the affine Schubert classes.\footnote{The geometric construction of this flat family uses the解析 geometry of the affine flag manifold; see \cite[Section 4]{Gr} for the technical details.}

This degeneration has the following key properties:
\begin{itemize}
\item It preserves the structure of the Schubert basis
\item It maps the equivariant structure constants $c_{u,v}^w(t,s)$ continuously to the affine counterparts
\item Since the degeneration limit is well-behaved, the positivity property of the coefficients is preserved
\end{itemize}

\section{Main Steps of the Proof}

Graham's proof can be divided into the following main steps:

\subsection{Step 1: Geometric Set-up}

Let $G = \mathrm{GL}_n(\mathbb{C})$, and let $B$, $B^-$, $T$ denote the standard Borel subgroup, the opposite Borel subgroup, and the maximal torus, respectively. We consider the corresponding affine geometry:
\begin{itemize}
\item The flag manifold $Fl_n = G/B$
\item The affine Grassmannian $Gr = G(\!(t)\!)/G[[t]]$
\item The affine flag manifold $\mathcal{F}\ell = G(\!(t)\!)/B(\!(t)\!)$
\end{itemize}

These spaces are naturally related, forming the geometric framework of the proof.

\subsection{Step 2: Construction of the Peterson Ring}

Define the Peterson ring $\mathcal{P}$ as
$$\mathcal{P} = H^*(Gr) \otimes \mathbb{Q}[t_1, \ldots, t_n] / (q_1, \ldots, q_{n-1}),$$
where $H^*(Gr)$ is the cohomology of the affine Grassmannian.

The Peterson isomorphism gives:
$$QH^*(Fl_n) \cong \mathcal{P}.$$

This ring possesses rich algebraic structure; its Schubert basis coefficients satisfy specific positivity properties.\footnote{The algebraic structure of $\mathcal{P}$ is closely related to the representation theory of affine $\mathrm{GL}_n$; see \cite{Gr} Section 3.}

\subsection{Step 3: Equivariant Extension}

To handle the equivariant case, Graham constructed an equivariant extension:
$$\mathcal{P}_T = H_T^*(Gr) \otimes \mathbb{Q}[t_1, \ldots, t_n] / (q_1, \ldots, q_{n-1}),$$
where $H_T^*(Gr)$ is the equivariant cohomology of the affine Grassmannian.

The key lemma (see \cite[Section 3]{Gr}) states that the Schubert basis elements in $\mathcal{P}_T$ satisfy the same positivity property as in the classical case.

\subsection{Step 4: Transversality and Degeneration}

The final step establishes the degeneration from equivariant Schubert classes to affine Schubert classes. Consider the family with parameter $\lambda$:
$$X_\lambda = \{ (g, z) \in G/B \times \mathbb{P}^1 : z \in \mathbb{P}^1 \},$$
as $\lambda \to \infty$, this family approaches an affine configuration.\footnote{The geometric meaning of this degeneration is related to the Tamagawa numbers of the group; see \cite{Gr} Section 5.}

The transversality in this degeneration process ensures the continuity of the coefficients, thereby transferring the positivity from the affine Grassmannian to the flag manifold.

\section{Relationship with Classical Schubert Calculus}

\begin{sloppypar}
An important corollary of Graham's theorem is that when restricted to the non-equivariant setting ($T$ acting trivially), the theorem reduces to the classical result.
\end{sloppypar}

\subsection{Non-Equivariant Degeneration}

When $T$ is the trivial torus (i.e., when we consider only ordinary cohomology), the coefficients $c_{u,v}^w(t,s)$ in \cref{th:GrahamMainPositivity} become ordinary integers $c_{u,v}^w$, and classical Schubert calculus tells us that $c_{u,v}^w \in \mathbb{Z}_{\geq 0}$.

This degeneration relationship can be expressed as:
$$c_{u,v}^w|_{t_i = s_j = 0} = c_{u,v}^w \in \mathbb{Z}_{\geq 0}.$$

\subsection{Connection with the Littlewood-Richardson Rule}

When both $u$ and $v$ are Grassmannian permutations, the Schubert classes can be identified with Schur functions. In this case, $c_{u,v}^w$ coincides exactly with the Littlewood-Richardson coefficients, which admit a combinatorial interpretation (counting semistandard Young tableaux).

Graham's positivity result can be viewed as a generalization of the Littlewood-Richardson coefficients to the equivariant setting.\footnote{The combinatorial interpretation in the equivariant case involves weighted tableaux; see \cite{Mihalcea} for further discussion.}

\section{Subsequent Developments and Influence}

Graham's work opened a new era in equivariant Schubert calculus.

\subsection{Mihalcea's Quantum Generalization}

In 2007, Leonardo Mihalcea extended Graham's result to the quantum setting in \cite{Mihalcea}, proving the positivity property in equivariant quantum Schubert calculus. This generalization has important applications in the Gromov-Witten invariant theory of flag manifolds.

Mihalcea's core idea is that although the quantum case is more intricate, one can reduce the problem to Graham's classical case through appropriate degeneration techniques.\footnote{The degeneration in the quantum case requires a careful analysis of the virtual structure; see \cite{Mihalcea} for details.}

\subsection{Connections with Representation Theory}

Graham's work reveals deep connections between Schubert calculus and representation theory. The cohomology of the affine Grassmannian is closely related to representations of affine Lie algebras, which has opened avenues for further interdisciplinary research.\footnote{These connections are further explored in the context of the geometric Langlands program; see \cite{Gr} for references.}

\section{Technical Details}

\subsection{The Affine Weyl Group}

The affine Weyl group $\widetilde{W}$ is the semidirect product of the Weyl group $W$ with the lattice $\mathbb{Z}^n$:
$$\widetilde{W} = W \ltimes \mathbb{Z}^n.$$
In the case of $\mathrm{GL}_n$, $\widetilde{W}$ is isomorphic to the double coset $S_n \ltimes \mathbb{Z}^n$.

Affine Schubert varieties are indexed by elements of the affine Weyl group and possess combinatorial properties similar to those of ordinary Schubert varieties.\footnote{The combinatorics of affine Schubert varieties is related to the theory of affine root systems; see \cite{Gr} Appendix A.}

\subsection{Moment Map and Transversality}

Let $K = U(n)$ be the compact subgroup of $G$, and let $T \subset K$ be the maximal torus. The moment map
$$\mu: X \to \mathfrak{t}^*$$
maps the flag manifold to the dual of the Lie algebra. Graham proved that the images of affine Schubert varieties under the moment map have good properties, which allows one to establish transversality results.\footnote{The properties of the moment map image are related to the theory of symplectic reduction; see \cite{Gr} Section 4.}

\section{Conclusion}

Graham's paper \cite{Gr} established a complete theory of positivity in equivariant Schubert calculus. Through the introduction of affine Schubert varieties and Peterson degeneration techniques, he proved the Peterson conjecture: namely, that when $u, v$ are Grassmannian permutations, the equivariant cup product coefficients $c_{u,v}^w(t,s)$ are polynomials with nonnegative coefficients in $t_i - s_j$.

This result not only resolved a long-standing conjecture but also laid the foundation for subsequent work in equivariant quantum Schubert calculus (Mihalcea) and more general positivity research. Graham's work has become a classic reference in Schubert calculus, widely cited and applied across the field.

\begin{Remark}
This chapter is based on reading \cite{Gr}. For questions or further discussion, please refer to the QA section in the appendix.
\end{Remark}
